\section{Implementation}

This chapter outlines the implementation approach for Med Flow. It covers the development process, core components, and technologies used. The system follows a component-based architecture with a FastAPI backend, MySQL and MongoDB databases, Faster-Whisper for transcription, Ollama or vLLM for clinical summarization, ReportLab for PDF generation, and SMTP for email delivery. The application is containerized with Docker for consistent deployment across environments. Configuration and logging are managed with Pydantic and the project's logging infrastructure.

\subsection{API Layer and Routing}

The API layer is built using FastAPI. It defines and serves HTTP endpoints and acts as the main entry point for user interactions.

The main API gateway is defined in \texttt{router.py}. This file creates the central API router and includes sub-routers for each domain, such as authentication, patients, consultations, transcriptions, review, prescriptions, PDF, email, audit, and administration. Each sub-router is responsible for a specific set of endpoints, which keeps the system separated by concern and easier to extend.

FastAPI dependency injection is used to manage authentication, database sessions, and shared resources. Dependencies are defined in \texttt{deps.py} and injected into endpoint functions. Each domain has its own router module where endpoints are defined as functions using decorators such as \texttt{@router.get} and \texttt{@router.post}.

\subsection{Authentication and Authorization}

Authentication ensures secure access control using role-based authorization. The system uses FastAPI dependency injection, JWT tokens, and a modular service structure.

The main authentication logic is implemented in \texttt{AuthApplicationService}, which bridges the API layer and domain logic. Domain models represent users, login requests, and the authentication interface, creating a consistent contract between application and infrastructure layers.

\begin{itemize}
  \item \textbf{User:} represents a system user with roles.
  \item \textbf{LoginRequest:} encapsulates login credentials.
  \item \textbf{AuthService:} defines an abstract interface for authentication logic and allows alternative implementations.
\end{itemize}

After successful authentication, the system assigns a JWT to the user. The token is used for later authenticated requests and includes user identity and role. Authentication checks are enforced through FastAPI dependencies that validate JWTs from incoming requests and ensure that only authenticated users can reach protected endpoints.

\subsection{Transcription Module}

The transcription component enables live conversion of audio into text, which can then be processed into clinical documentation. It is designed to be scalable, replaceable, and integrated with the other components.

The workflow is orchestrated by \texttt{TranscriptionApplicationService}, which handles streaming audio data, temporary audio and text chunks, and finalization of the complete transcription. The actual speech-to-text conversion is handled by Faster-Whisper through a dedicated adapter. The adapter communicates with a local Whisper microservice using HTTP.

\begin{enumerate}
  \item \textbf{Session initialization:} temporary storage is created for incoming audio data through \texttt{start\_transcription\_session}.
  \item \textbf{Audio streaming and chunk handling:} audio is streamed to the backend. The service processes chunks and sends them to the Faster-Whisper API for transcription. Partial results are stored using the temporary transcript chunk repository.
  \item \textbf{Transcript assembly:} temporary chunks are assembled into the final transcription, which is then stored in MongoDB as a consultation document.
  \item \textbf{Error handling:} the module handles network interruptions, partial uploads, and API errors. Temporary storage ensures that progress is not lost if a session is interrupted.
\end{enumerate}

\subsection{LLM Orchestration}

The LLM orchestration component enables clinical note creation, report generation, and suggestions using the locally integrated Ollama model. The core logic lives in LLM client and adapter classes, such as \texttt{ollama\_client.py} and \texttt{ollama\_adapter.py}. These classes provide an interface for sending prompts, receiving completions, validating structured output, and handling model health checks.

The \texttt{OllamaClient} handles HTTP communication with the local Ollama or vLLM runtime. Application services or endpoints use predefined prompt templates stored in the prompt directory. These templates are populated with clinical context, patient data, transcript content, and other inputs before being sent to the model runtime.

The orchestration workflow is:

\begin{enumerate}
  \item \textbf{Prompt construction:} templates are populated with patient context and consultation content.
  \item \textbf{Request dispatch:} the prompt is sent through \texttt{OllamaClient}, which handles retries, timeouts, and errors.
  \item \textbf{Response handling:} the response is returned to the calling service and can be displayed directly or processed further by Suggestive Mode.
  \item \textbf{Error handling:} network errors, model unavailability, and malformed responses are handled by the orchestration layer to provide informative errors.
\end{enumerate}

\subsection{PDF Generation}

The PDF generation module creates structured PDF documents from clinical data produced by consultations and prescriptions. This feature is important for shareable and printable records.

PDF generation is implemented with ReportLab. An adapter class encapsulates formatting and creation logic, while the service orchestrates data retrieval, PDF creation, and storage. The workflow starts with collecting consultation, patient, staff, and prescription data from repositories. The \texttt{ReportLabPdfGenerator} then formats the data into a document with page layouts, styles, tables, and visual elements. The resulting file is stored and made available for download or email attachment after approval.

\subsection{Email Notification Service}

The email notification service sends emails to patients containing consultation summaries and prescription information. It supports communication between the application and its users after doctor approval.

Email sending is handled by \texttt{SmtpEmailService}. It constructs and sends emails using the SMTP library \texttt{smtplib}. The service supports plain text and HTML emails, as well as attachments such as PDF files. Multipart messages are built with Python's \texttt{email.mime} package.

During development, MailHog was used as the SMTP server. MailHog captures outgoing emails so developers can view and verify email content without sending real messages. The email workflow constructs an email message, connects to the SMTP server, sends the email, and then allows developers to inspect the captured message through MailHog's web interface.

\subsection{Database Integration}

The system uses a dual-database architecture to manage structured and unstructured data. MySQL is used for structured data such as user accounts, patient records, consultations, prescriptions, and audit logs. MongoDB is used for semi-structured data such as clinical documents, consultation transcripts, generated reports, and finalized AI artifacts.

SQLAlchemy manages the MySQL connection and session lifecycle. Reusable database access is provided through functions such as \texttt{get\_engine()} and \texttt{get\_session()}. PyMongo manages MongoDB connections, with \texttt{get\_database()} returning the database handle used by document repositories.

Both databases are abstracted behind repository classes. Application services interact with repositories rather than directly with database clients. This pattern improves separation of concerns, testability, and maintainability.